\documentclass[12pt, twoside]{article}
\input{../preamble}

\fancyhead[LO]{La Scuola d'Italia / Dr.\ Huson / IB Math \\* 6 May 2026}
\fancyhead[RO]{\fbox{\begin{minipage}{6cm}\footnotesize First \& Last name:\\[0.5cm]\end{minipage}}}
\fancyfoot[C]{\thepage}

\begin{document}
\subsubsection*{7.6 Classwork: Natural Logarithms and $\exp$ (no calculator)}

\begin{enumerate}[itemsep=0.4cm]

%--- Problem 1: Warmup --- evaluate logs (integers + variables) [8 marks] ---
\item[1.] \textbf{Warmup}

\textbf{(a)} Evaluate each of the following.

\begin{enumerate}[label=(\roman*), itemsep=1cm]
  \item $\log_9 27$
  \vspace{0.6cm}

  \item $\log_{27} 9$
  \vspace{0.6cm}

  \item $\log_{1/3} 27$
  \vspace{0.6cm}

  \item $\log_2 \sqrt[4]{8}$
  \vspace{1cm}
\end{enumerate}

\textbf{(b)} Let $a > 1$ and $b > 0$. Write each expression in terms of
$\log_a b$.

\begin{enumerate}[label=(\roman*), itemsep=1cm]
  \item $\log_a (a^3 b)$
  \vspace{0.6cm}

  \item $\log_a \!\left(\dfrac{b^2}{a}\right)$
  \vspace{0.6cm}

  \item $\log_a \sqrt{a^3\, b}$
  \vspace{0.6cm}

  \item $\log_{a^3} b$
  \vspace{0.6cm}
\end{enumerate}

\newpage

%--- Problem 2: Simplify expressions with e and ln (exact) [8 marks] ---
\item[2.] Simplify each of the following. Give exact answers.

\begin{enumerate}[label=(\alph*), itemsep=0.6cm]
  \item $\ln\!\left(e^{\,4}\right)$
  \vspace{0.9cm}

  \item $e^{\,\ln 7}$
  \vspace{0.9cm}

  \item $\ln\!\left(e^{\,2x+1}\right)$
  \vspace{0.9cm}

  \item $e^{\,3\ln 2}$
  \vspace{0.9cm}
\end{enumerate}

%--- Problem 3: Solve simple e^x and ln x equations [8 marks] ---
\item[3.] Solve each equation for $x$. Give exact answers in terms of $e$ or $\ln$.

\begin{enumerate}[label=(\alph*), itemsep=0.6cm]
  \item $e^{\,x} = 12$
  \vspace{0.9cm}

  \item $\ln x = 5$
  \vspace{0.9cm}

  \item $e^{\,2x-1} = 9$
  \vspace{0.9cm}

  \item $\ln(3x + 1) = 2$
  \vspace{0.9cm}
\end{enumerate}

\newpage

%--- Problem 4: Exponential growth (bacteria, continuous model) [13 marks] ---
\begin{center}
\textbf{A graphing calculator is allowed for the remaining problems.}
\end{center}
\vspace{0.2cm}

\item[4.] A scientist places $250$ bacteria in a Petri dish. The number of
bacteria $N$, $t$ hours after the start of the experiment, is modelled by
$$N(t) \;=\; 250\, e^{\,k t}, \qquad k > 0.$$
After $3$ hours, the dish contains $1000$ bacteria.

\begin{enumerate}[label=(\alph*), itemsep=0.3cm]
  \item Write down the value of $N$ when $t = 0$.

  \vspace{1.6cm}

  \item Show that $k = \dfrac{1}{3}\ln 4$.

  \vspace{2.6cm}

  \item Find the number of bacteria after $5$ hours. Give your answer to
  the nearest integer.

  \vspace{2.4cm}

  \item Find the time, in hours, at which the number of bacteria first
  reaches $10\,000$. Give your answer correct to three significant
  figures.

  \vspace{2.4cm}

  \item Find the doubling time of the population. Give your answer
  correct to three significant figures.

  \vspace{2cm}
\end{enumerate}

\newpage

%--- Problem 5: Continuous vs annual compound interest [13 marks] ---
\item[5.] Lucia invests \$5000 in an account that earns interest at a
nominal annual rate of $6\,\%$, compounded continuously. The value of the
account, in dollars, after $t$ years is given by
$$A(t) \;=\; 5000\, e^{\,0.06\, t}.$$

\begin{enumerate}[label=(\alph*), itemsep=0.3cm]
  \item Find the value of the account after $4$ years. Give your answer
  correct to the nearest cent.

  \vspace{2.5cm}

  \item Find the time required for the value of the account to first
  reach \$8000. Give your answer correct to three significant figures.

  \vspace{3cm}

  \item Find the time required for the value of the account to double.
  Give your answer correct to three significant figures.

  \vspace{3cm}

  \item A second account also begins with \$5000 but earns $6\,\%$
  compounded annually. Find the value of this second account after
  $4$ years, and state the difference between the two accounts at that
  time. Give both answers correct to the nearest cent.

  \vspace{4cm}
\end{enumerate}

\end{enumerate}
\end{document}
